\section{The inherited primitive fiber interface}
\label{sec:tpc21-interface}

Fix \(h\in\Z\setminus\{0\}\), put \(H=\rad|h|\), and let
\[
 e(t)=\exp(2\pi i t),\qquad e_s(t)=e(t/s).
\]
\begin{equation}
 R=\lfloor X^{1/2-\delta}\rfloor,
 \qquad
 V=\lfloor R^{1/2}\rfloor,
 \qquad 0<\delta<\frac12.
 \label{eq:tpc21-RV}
\end{equation}
Take \(X\) large enough that \(R\ge H^2\).  Then the inherited
restriction \(d\le V\) implies \(dH\le R\), as required in the TPC-20
recombination.
The row parameters satisfy
\begin{equation}
 Q=LD,\qquad QJ\asymp X,\qquad
 L>\max(2R,2|h|),
 \label{eq:tpc21-row-scales}
\end{equation}
and
\begin{equation}
 \ell\in[L/2,2L]\text{ is prime},
 \qquad
 d\in[D,2D]\cap[1,V],
 \qquad
 m=\ell d.
 \label{eq:tpc21-row-support}
\end{equation}
The integer \(d\) is squarefree, \((d,H)=1\), and the orbit is restricted
to \((j,H)=1\).  In particular \((m,H)=1\).  All implied constants may
depend on \(h,\delta\), and fixed smooth weights.

TPC-20 supplies two short--short coefficient choices,
\begin{equation}
 \vartheta_R(u)=\lambda'_R(u)
 \quad\text{or}\quad
 \vartheta_R(u)=b_R(u)\one_{u\le R},
 \label{eq:tpc21-theta-choices}
\end{equation}
where both are supported on squarefree \(u\le R\) and, for every
fixed \(\eps>0\),
\begin{equation}
 |\vartheta_R(u)|\ll_\eps X^\eps.
 \label{eq:tpc21-theta-soft}
\end{equation}
Both coefficient choices are real-valued.  This fact is used whenever
a Hermitian packet is written as an absolute square.
The first choice is the centered finite Ramanujan model.  The second is
the matched residual coefficient restricted to its two short divisor
legs.  No statement below controls the complementary terms with
\(u>R\) or \(v>R\).

\subsection{Rows, masks, and fiber vectors}

Let \(\cA\) be one finite separated row packet.  Its labels
\(\alpha=(\ell,d)\) satisfy \eqref{eq:tpc21-row-support}, and its
weighted cardinality obeys
\begin{equation}
 \sum_{\alpha\in\cA}|x_\alpha|
 +\sum_{\alpha\in\cA}|y_\alpha|
 \ll QX^\eps
 \label{eq:tpc21-row-mass}
\end{equation}
for the bounded row weights used below.  The original smooth factors in
\(mj/X\) and \(dj/K\) admit the standard rapidly convergent Mellin--Fourier
separation.  Each separated term has a common \(F(j/J)\) and rank-one
row weights \(x_{\alpha_1}y_{\alpha_2}\); the total projective mass of
the truncated superposition is \(X^\eps\), with arbitrary power error.

The generic TPC-18 row-pair mask is
\begin{equation}
 \mathfrak m(\alpha_1,\alpha_2)
 =\one_{\{\ell_1\ne\ell_2\}}
  \one_{\{|m_1-m_2|>QX^{-\kappa}\}}
  \one_{\{(d_1,d_2)\le X^\kappa\}},
 \label{eq:tpc21-generic-mask}
\end{equation}
where \(\kappa>0\) is fixed.  We also allow an arbitrary fixed
row-pair restriction independent of \(j\), denoted by the same letter,
when proving the exact fiber identities.

For squarefree \(u,v\le R\), put
\begin{equation}
 g=(u,v),\qquad q=[u,v].
 \label{eq:tpc21-gq}
\end{equation}
All aggregated layers used below are restricted to \((q,H)=1\).
Layers with \((q,H)>1\) have empty admissible fibers and are omitted;
in particular \(\overline H\pmod q\) is always defined.
If
\[
 (u,m_1H)=(v,m_2H)=1,
\]
then the pair \((m_1,m_2)\) is compatible precisely when
\(g\mid m_1-m_2\).  It then determines the unique unit class
\begin{equation}
 \kappa\in\cU_q=(\Z/q\Z)^\times,
 \qquad
 m_1\kappa\equiv-h\pmod u,
 \quad
 m_2\kappa\equiv-h\pmod v.
 \label{eq:tpc21-joint-kappa}
\end{equation}
If either row--modulus admissibility condition fails, the corresponding
fiber is zero.
Define
\begin{align}
 Z_{u,v}^{\mathfrak m}(\kappa)
 ={}&\sum_{\alpha_1,\alpha_2\in\cA}
 x_{\alpha_1}y_{\alpha_2}
 \mathfrak m(\alpha_1,\alpha_2)
 \one_{\{m_1\equiv-h\overline\kappa\ (u)\}}
 \one_{\{m_2\equiv-h\overline\kappa\ (v)\}},
 \label{eq:tpc21-uv-fiber}\\
 Z_q^{\mathfrak m}(\kappa)
 ={}&\sum_{\substack{u,v\le R\\{}[u,v]=q}}
 \vartheta_R(u)\vartheta_R(v)
 Z_{u,v}^{\mathfrak m}(\kappa).
 \label{eq:tpc21-q-fiber}
\end{align}
Congruences modulo \(1\) are vacuous.  An inadmissible divisor--row
combination contributes zero.  In the retained low nonzero joint
frequencies of TPC-20, the common multiplier is \(g\), and the factor
\(g/(uv)=1/q\) gives the scalar packet
\begin{equation}
 \cL_q
 =\frac{J}{Hq}\sum_{r\ne0}c_H(r)
 \widehat F\!\left(\frac{rJ}{Hq}\right)
 S_q(r),
 \qquad
 S_q(r)=\sum_{\kappa\in\cU_q}
 Z_q^{\mathfrak m}(\kappa)e_q(r\overline H\kappa).
 \label{eq:tpc21-long-packet}
\end{equation}

Put
\begin{equation}
 \overline Z_q^{\mathfrak m}
 =\frac1{\varphi(q)}\sum_{\kappa\in\cU_q}
 Z_q^{\mathfrak m}(\kappa),
 \qquad
 D_q^{\mathfrak m}(\kappa)
 =Z_q^{\mathfrak m}(\kappa)-\overline Z_q^{\mathfrak m}.
 \label{eq:tpc21-mean-discrepancy}
\end{equation}
Then
\begin{equation}
 S_q(r)=\overline Z_q^{\mathfrak m}c_q(r)
 +\widehat D_q^{\mathfrak m}(r),
 \qquad
 \sum_{r\bmod q}|\widehat D_q^{\mathfrak m}(r)|^2
 =q\lVert D_q^{\mathfrak m}\rVert_{2,\cU_q}^2.
 \label{eq:tpc21-additive-parseval}
\end{equation}
The mean term is closed in \cref{sec:tpc21-mean-closure}; the discrepancy is
the subject of the remaining sections.

\subsection{The inherited restricted-frequency gate}

For a long-modulus family \(\cQ\), TPC-20 proves, after truncating the
rapidly decaying Fourier tail, that
\begin{equation}
 \sum_{q\in\cQ}|\cL_q^{\rm disc}|
 \ll_{H,F,\eps}X^{\eps/2}\sqrt{\frac JH}
 \sum_{q\in\cQ}\lVert D_q^{\mathfrak m}\rVert_{2,\cU_q}
 +O_A(X^{-A}\cZ_1),
 \label{eq:tpc21-inherited-gate}
\end{equation}
under its no-wrap and polynomial-mass hypotheses.  Here
\(\cZ_1\) denotes the relevant total coefficient mass.  The natural
row-pair--orbit scale is
\begin{equation}
 Q^2J\asymp XQ.
 \label{eq:tpc21-natural-scale}
\end{equation}
Every use of \eqref{eq:tpc21-inherited-gate} below is compared with
this scale.  The estimate is sufficient, not known to be necessary;
in particular, scalar cancellation may be invisible to the nonlinear
fiber norm.
